\section{Discussion}

We have argued that a transdiagnostic dimensional map is most useful when read not
as a table of coordinates but as a measurement instrument, and we have shown that a
single missingness-native Bayesian bifactor model makes five instrument properties
computable and testable on one object: it projects each patient to a coordinate
with its own uncertainty, weighs each indicator by the information it contributes,
prescribes how few items suffice and where the item bank itself is the limit,
values the cheapest next measurement, and --- as we verified by simulation ---
reports uncertainty that is calibrated and a low-dimensional summary whose fidelity
is quantified. None of the individual components is new: expected-a-posteriori
scoring, Fisher information, adaptive testing, value-of-information design,
full-information treatment of missing data, and archetypal analysis are each
standard. The contribution is their integration, and their joint validation on a
real, pervasively incomplete, mixed-type transdiagnostic assessment at $N=9{,}013$
--- the setting where each component is usually assumed rather than demonstrated.

The instrument framing changes what a map is for. A point-estimate map invites the
reader to treat every coordinate as equally trustworthy and every axis as equally
measured; the honest answer, visible in Figure~\ref{fig:projection}, is that a
patient may be pinned to a posterior standard deviation of $0.22$ on one axis and
floated in from the prior on another, and that the map should say so. Once the map
reports information rather than only position, three practical capabilities follow
directly. First, the item bank can be audited: Figure~\ref{fig:information} shows
that loading and information part company, so a battery assembled to maximise
loadings can be dominated by rare flags that localize no one, and the fix ---
reorder by information --- is immediate. Second, assessment can be made cheaper
without guessing: the minimal-indicator rule and the value-of-information battery
(Figs.~\ref{fig:adaptive},~\ref{fig:voi}) turn ``which items do we actually need''
into a closed-form calculation, and the finding that a fixed most-informative-first
order is a tight upper bound on adaptive efficiency means the saving can be shipped
as a static list rather than requiring adaptive-testing infrastructure. Third,
under-measured cohorts are localized rather than assumed: the same accounting shows
exactly which axis in which cohort would benefit most from additional measurement.

These capabilities would be hollow if the instrument's machinery did not beat the
simpler tools a reader already has, so we tested it against them directly. Ranking
battery items by information rather than loading is not a cosmetic preference: at a
fixed $27$-item budget it buys $0.11$ of mean reliability
(Fig.~\ref{fig:added}A), because loading and prevalence pull in different
directions. Carrying the full posterior coordinate rather than a per-axis sum-score
improves out-of-sample prognosis, and does so most where it matters ---
under sparse measurement, the regime in which a point estimate is least trustworthy
(Fig.~\ref{fig:blockers}A). And the calibration that a same-model simulation can only
show to be internally consistent survives genuine misspecification: coverage degrades
smoothly, never catastrophically, when the data-generating process breaks local
independence, omits a factor, or violates Gaussianity
(Fig.~\ref{fig:blockers}B). These are the comparisons a measurement venue asks for. On the two design
decisions the instrument is put to --- which items to collect, and which coordinate
to carry downstream --- it wins outright, not by a wide margin on any single one but
consistently, and for mechanistic reasons the framework predicts in advance. On the
third, reconstruction fidelity, it does not win and does not need to: we show the
metric itself cannot separate a convex summary from a random subspace on this
near-isotropic cloud, so the archetype simplex earns its place on interpretability
and graded named membership rather than on variance retained. Winning where the
comparison is decision-relevant, and dismantling the yardstick where it is not, is
the honest reading of the four analyses.

Two results are worth stating as limits the instrument imposes on itself. The
bank-limited axes --- mania/activation and substance, capped at reliability $0.41$
and $0.43$ --- cannot be improved by any amount of adaptive scheduling or data
completion, because the information is not present in the items; the same two axes
are the ones the archetype summary reconstructs worst
(Fig.~\ref{fig:archetype}C). This is not a failure of the method but a diagnosis it
delivers: these constructs need better instruments, not better analysis. And the
five-corner convex summary retains $59\%$ of coordinate variance against an $80\%$
linear optimum --- but the honest reading of that gap is not that the simplex is a
worse projection. The map's coordinate variance is nearly isotropic, so a
\emph{random} five-dimensional subspace already retains $62\%$ (the analytic $k/d$),
and the archetype simplex sits slightly below it: reconstruction $R^2$ is simply the
wrong axis on which to judge this summary, because a coin-flip subspace beats it and
a principal component --- a signed contrast no patient can \emph{be} --- beats it
further while providing nothing a clinician can act on. The simplex earns its place
not by variance retained but by what it uniquely supplies: convex, named, per-patient
membership in five clinical poles over a population that has no natural clusters to
begin with (silhouette $z=1.13$). Stating this plainly, rather than quoting the $R^2$
one would prefer, is itself the point: an instrument that reports the axis on which
it should \emph{not} be judged is more useful than one that claims a fidelity victory
it does not have.

Several limitations bound the claims. Calibration is established in simulation
under the model's own generative assumptions and real missingness patterns; it
demonstrates that the projection machinery is internally honest --- that its
$95\%$ intervals cover at $95\%$ --- but not that the fitted model is the true data
generating process, which no coverage study can. Posterior standard deviations run
slightly below residual error on the discrete-informed axes (a light-tailed Laplace
approximation), so interval width is marginally optimistic even though the
$95\%$-level coverage, set by the interval quantiles, is exact. The
value-of-information accounting is greedy and Gaussian-entropy based, an upper bound
on a jointly-optimal design. The reliabilities and item counts are properties of
this bank in these cohorts and will differ for another assessment --- which is
precisely why the instrument is designed to be re-run on a new one. Finally, the
adaptive-testing analysis is a single-axis simulation on model-generated responses;
extending it to the joint eight-axis posterior with cross-axis information sharing
is a natural next step that would only tighten the case for a fixed battery.

This paper deliberately reports no clinical or biological finding. The immunometabolic
axis that recurs throughout as a worked example, the archetype corners, the
cohort-level prognosis --- these are substantive results, but they belong to the
companion paper, where they are established on their own terms with the appropriate
clinical framing and controls. The division is not cosmetic: a reader who wants to
know whether the immunometabolic axis predicts remission should not have to
adjudicate a measurement-model derivation to get there, and a reader who wants to
reuse the instrument on a new cohort should not have to extract it from a clinical
narrative. Kept separate, each paper is self-contained; the measurement instrument
described here is what makes the clinical map trustworthy, and the clinical map is
what makes the instrument worth building.
